\subsection{Chromosome-arm profiles and sequencing-run consistency}\label{sec:m-y-followup}
Arm-level quality control covered all 1,639 processed samples. Single-copy Y marker counts were summed separately over GRCh38 Yp (positions below 10.3 Mb) and proximal Yq (10.6--26.6 Mb) and divided by chromosome 2 marker counts. Each arm ratio was scaled to its median across samples with both raw ratios above half their initial cohort medians. These proximal-Yq markers do not measure the distal Yq12 satellite block.

The priority-ordered screening rule assigned negligible Y signal when whole-Y dosage was at most 0.0222, the highest value in 65 held-out female genomes; Yp depletion when Yp was below 0.5 and Yq within 0.1 of one; Yq loss when Yp was within 0.1 of one and Yq more than 0.1 below Yp; and arm imbalance when the absolute Yp--Yq difference exceeded 0.1. Among remaining samples, whole-Y dosage of at least 1.75, 1.25--1.75 or below 0.9 flagged Y gain, intermediate gain or low dosage, respectively. The remaining samples had no flag. The corresponding labels (\texttt{no\_Y}, \texttt{yp\_deletion}, \texttt{yq\_loss}, \texttt{partial}, \texttt{Y\_gain}, \texttt{mosaic\_gain}, \texttt{mosaic\_loss} and \texttt{normal}) are screening categories, not confirmed karyotypes, deletion boundaries or mosaic fractions. Counts were 42 (41 females and NA19226), 6, 1, 3, 2, 3, 11 and 1,571, respectively. Y dosage is measured against chromosome 2, so it does not name a karyotype: a separate column combines the X and Y dosages and gives XY (1,577), XX (39 of the 41 females), XXY (1), XYY (2), X (1) or, when either dosage is more than 0.15 from a whole number, mosaic (19: 17 males, including the samples with intermediate Y dosage, and two females with X dosage 1.795 and 1.848).

Six samples showed Yp ratios of 0.260--0.277 with proximal-Yq ratios of 0.985--1.033 (Table~\ref{tab:y-arm-flags}). They include father--son pairs HG01982--HG01984 and HG03006--HG03008, consistent with inherited Yp variation. Recurrent multigene Yp deletions involving TSPY have been described \cite{Jobling2007}; sequence confirmation is needed to identify the specific deletion. In unflagged samples, Yp contributes a median 12.55\% of ruler hits, versus 3.63--3.75\% in these six. Under stable proximal-Yq dosage, the calibrated arm-composition diagnostic predicts satellite-length inflation of 10.06--10.20\%. This replaces the earlier uncalibrated Yq-to-total-dosage ratio. We excluded the uncorrected estimates; a Yq-only denominator requires separate calibration.

Three further samples had Yp/Yq imbalance (HG01967, HG02966 and NA20806) and one had Yq loss with an intact Yp (HG00185). Together with the six Yp-depleted samples, they are retained in displays and excluded from population summaries. Eleven samples with concordantly reduced Yp and Yq signal and estimable lengths are included despite dosage below 0.9; the lowest are HG00246 and NA20754 (Y dosage 0.3227 and 0.4130), which an earlier rule based on arm ratios below 0.5 had called negligible. A uniform multiplicative reduction in Y abundance cancels between the array and single-copy Y signals, whereas loss or gain of selected ruler regions need not cancel. This algebraic distinction motivates the inclusion rule; it does not establish the biological cause of low signal or identical measurement uncertainty at every Y dosage.

For the HG02965--HG02966 pair, reads were partitioned by flowcell using the CRAM read-group metadata and processed separately (Table~\ref{tab:y-flowcells}). Each sample has one recorded library spanning three flowcells, with four read groups per flowcell. Relative to the same individual's whole-library result, the son's Yp signal ranged from 0.991 to 1.019, proximal-Yq signal from 0.9996 to 1.0004, DYZ3 from 0.995 to 1.007 and the near-absent Yq12 signal from 0.964 to 1.053. The near-absent DYZ19 signal varied more (0.864--1.201). The whole-library arm ratios were 1.910 for Yp and 1.082 for proximal Yq in the son. Thus, Yp excess and distal satellite depletion recur across sequencing runs. This argues against a flowcell-specific anomaly; shared library preparation and cell-line history remain possible sources. These are technical partitions of one library, not independent DNA preparations or independent validation of chromosome structure.

Nanopore sequencing of HG01967 (ENA ERR15140186, downloaded with checksum verification) also showed increased Y signal and reduced Yq12 abundance: dosage was 1.389 and the raw Yq12-to-ruler count ratio was 0.782. The corresponding short-read quantities were approximately 1.35 and 0.625. Agreement in direction across platforms supports the unusual profile. The marker-count ratios are not independently calibrated long-read array lengths, and these data do not establish rearrangement boundaries.

Profiles along CHM13v2 chrY place these signals (Supplementary Fig.~\ref{fig:chrY-unusual}). Each 100 kb window shows, relative to the median of 100 XY males with Y state normal, the MQ20 read depth moved from GRCh38 with the UCSC hg38ToHs1 chain where reads map uniquely, and elsewhere the median count of the yruler marker k-mers that match the window exactly; 62 of 622 windows, all in segmental duplications, have neither. In HG02966, chrY is at two copies from Yp to 17.10 Mb and absent from 17.30 Mb through Yq12, and the centromere window is at 2.04 copies. The transition lies between the two arms of palindrome P6 (chrY:17,066,095--17,176,113 and 17,222,323--17,332,346) \cite{Rhie2023}, the configuration expected for an isodicentric Y formed by recombination between palindrome arms \cite{Lange2009}; his father HG02965 shows no change. HG01967 shows the same pattern at lower amplitude (1.68 copies in Yp11.2 and q11.221, 0.30 in q11.223 and 0.27 in Yq12, consistent with 68\% of cells carrying such a chromosome), with the transition in q11.222. In 15 father--son pairs with an unusual member (Supplementary Fig.~\ref{fig:chrY-pairs}), Yp depletion (two pairs), a long DYZ18S array and a short DYZ19 array were shared by father and son, whereas mosaic gains and losses, the Y gain of HG03456, the XXY state of HG01683, the HG02966 rearrangement and the Y loss of NA19226 were each seen in one generation. Read depth does not distinguish germline changes from changes that arose in the cell lines.

\begin{table}[p]
\caption{Samples with arm-imbalanced or Yp-depleted marker profiles. Arm ratios are cohort-normalized; whole-Y dosage comes from the current population estimate table. All ten samples are excluded from population summaries. These profiles require structural follow-up.}\label{tab:y-arm-flags}
\centering\small
\begin{tabular}{llrrr}
\toprule
Sample & Profile & Y dosage & Yp & Proximal Yq \\
\midrule
HG00185 & Yq loss & 0.813 & 1.015 & 0.756 \\
HG01967 & Yp/Yq imbalance & 1.350 & 1.683 & 1.259 \\
HG01982 & Yp depletion & 0.923 & 0.260 & 0.990 \\
HG01984 & Yp depletion & 0.932 & 0.267 & 0.998 \\
HG02966 & Yp/Yq imbalance & 1.223 & 1.911 & 1.083 \\
HG03006 & Yp depletion & 0.921 & 0.267 & 0.985 \\
HG03008 & Yp depletion & 0.931 & 0.269 & 0.997 \\
HG03976 & Yp depletion & 0.934 & 0.272 & 0.999 \\
HG03998 & Yp depletion & 0.965 & 0.277 & 1.033 \\
NA20806 & Yp/Yq imbalance & 1.030 & 1.143 & 0.980 \\
\botrule
\end{tabular}
\end{table}

\begin{table}[p]
\caption{Technical consistency across flowcells for father HG02965 and son HG02966. Each sample has one recorded library. Depth is the mean chromosome 2 marker count. DYZ3, Yq12 and DYZ19 columns give ratios to the same sample's whole-library measurement; the son's near-absent satellite signals must not be interpreted as normal absolute lengths.}\label{tab:y-flowcells}
\centering\scriptsize
\begin{tabular}{llrrrrr}
\toprule
Sample & Flowcell & Depth & Y dosage & DYZ3 & Yq12 & DYZ19 \\
\midrule
HG02966 & HV3CYDSXX & 7.67 & 1.225 & 0.995 & 0.969 & 0.882 \\
HG02966 & HV3GMDSXX & 7.88 & 1.232 & 0.996 & 0.964 & 0.864 \\
HG02966 & HV3HYDSXX & 9.81 & 1.226 & 1.007 & 1.053 & 1.201 \\
HG02965 & HTGLNDSXX & 9.30 & 1.046 & 1.010 & 1.010 & 0.997 \\
HG02965 & HV3CGDSXX & 8.55 & 1.049 & 0.996 & 0.995 & 0.993 \\
HG02965 & HV3KJDSXX & 8.66 & 1.047 & 0.993 & 0.994 & 1.011 \\
\botrule
\end{tabular}
\end{table}
